\subsubsection{Analisis STRIDE pada Sistem DecMed}
\label{subsec:analisis-stride-decmed}

Analisis STRIDE digunakan untuk mengevaluasi implikasi keamanan dari arsitektur DecMed saat ini ketika ditempatkan pada skenario layanan kesehatan kolaboratif yang menjadi fokus analisis. Berdasarkan pembahasan pada Subbab \ref{sec:stride}, tahapan awal pemodelan ancaman menggunakan STRIDE adalah mengidentifikasi komponen dan alur kerja sistem yang menjadi target ancaman. Identifikasi tersebut dilakukan melalui penggambaran DFD yang menunjukkan aliran data dan komunikasi antarkomponen sistem. Pada analisis ini, DFD Level 0 dan Level 1 dimodelkan. Pemodelan ke level lebih lanjut tidak dilakukan karena pemodelan hingga level 1 sudah cukup untuk mengidentifikasi komponen dari DecMed dan alur data antarkomponen. Hasil pemodelan DFD dapat dilihat pada Lampiran \ref{lampiran:dfd-decmed}.

Berdasarkan arsitektur DecMed pada Gambar \ref{fig:decmed-archi} dan DFD pada Lampiran \ref{lampiran:dfd-decmed}, \textit{attack surface} utama sistem mencakup \textit{Patient Client}, \textit{Hospital Client}, \textit{Ministry Client}, data pada lapisan \textit{on-chain}, data terenkripsi pada lapisan \textit{off-chain}, Redis PRE, jaringan IOTA, IOTA \textit{Gas Station}, serta token kapabilitas yang mengatur otorisasi akses. Hasil identifikasi ancaman berdasarkan metode STRIDE terhadap komponen utama pada DecMed dapat dilihat pada Lampiran \ref{lampiran:analisis-stride-decmed}.

Berdasarkan hasil analisis STRIDE pada Tabel \ref{tab:stride-decmed}, ancaman yang paling relevan terhadap fokus analisis adalah \textit{spoofing} (T01), \textit{tampering} (T02 dan T03), \textit{information disclosure} (T05 dan T06), dan \textit{elevation of privilege} (T08). Ancaman-ancaman tersebut secara langsung berkaitan dengan kebutuhan akan model \textit{policy} akses yang lebih luas, segmentasi data \textit{off-chain}, dan mekanisme delegasi yang mampu membatasi hak akses secara bertahap. Ancaman \textit{repudiation} (T04) tetap relevan terhadap kebutuhan audit delegasi, tetapi diposisikan sebagai implikasi pendukung karena fokus utama berada pada perluasan \textit{policy} akses dan delegasi. Ancaman T03 dan T06 ditangani secara terbatas melalui kebutuhan nonfungsional \texttt{KNF-TA-07} dan \texttt{KNF-TA-08}, yaitu deteksi modifikasi \textit{ciphertext} \textit{off-chain} dan ketiadaan \textit{payload} klinis plaintext pada metadata \textit{on-chain}. Ancaman \textit{denial of service} (T07) tidak dijadikan masalah yang diselesaikan ataupun diuji, karena tidak termasuk dalam fokus terkait model \textit{policy} akses, segmentasi data, dan mekanisme delegasi antartenaga kesehatan.

Dengan demikian, hasil analisis STRIDE memperkuat bahwa keterbatasan pada model kontrol akses dan delegasi DecMed tidak hanya merupakan persoalan fungsional, tetapi juga memiliki implikasi keamanan yang nyata. Oleh karena itu, diperlukan mekanisme otorisasi yang mampu mendukung pembatasan akses yang lebih rinci dan aman pada sistem rekam medis elektronik terdesentralisasi